%%%%

%%%   Самарский государственный технический университет

%%%   Кафедра прикладной математики и информатики

%%%

%%%   Преамбула для статьи в журнал "Вестник Самарского государственного

%%%   технического университета. Серия: «Физико-математические науки»"

%%%   Ориентирована под MikTEx 2.4 for Win32

%%%

%%%   Хотя сам журнал собирается под GNU/Linux на дистрибутиве TeXLive



\documentclass[11pt,a4paper,twoside]{article}



\usepackage[cp1251]{inputenc}

\usepackage[T2A]{fontenc}

\usepackage[english,russian]{babel}

\usepackage{samgtubib}

\usepackage[dvips]{graphicx}

\usepackage[multiple]{footmisc}



\usepackage{samgtu}

\renewcommand{\VestnikYear}{2009} % Год издания Вестника

\renewcommand{\VestnikNumber}{2\,(19)} % Номер Вестника

\renewcommand{\VestnikMonth}{Ноябрь} % Месяц выхода Вестника

\renewcommand{\VestnikData}{01.11.09} % Дата подписания Вестника в печать

\renewcommand{\VestnikUPL}{26,64} % Усл. печ. л.

\renewcommand{\VestnikUIL}{25,73} % Уч.-изд. л.

\renewcommand{\VestnikRegNumber}{479/09} % Регистрационный номер

\renewcommand{\VestnikOrder}{994} % Номер заказа






%
%
%
%%\samgtupubl % для генерации pdf, отдаваемого в типографию СамГТУ
%\pdfcrop % обрезка pdf для размещения на math-net.ru
%\draft % На корректуре
\filllastpage % Последняя страница не заполнена не полностью
%\briefcomm % Статья является кратким сообщением

\vsgtu{1266}

\FirstPage{136}
\LastPage{148}
%
%
%\begin{document}

\renewcommand{\baselinestretch}{0.923}

\udk{004.032.26:004.827}
\msc{93B40, 42C05, 33C45}

\rutitle{Решение задачи классификации с использованием нейронных
неч\"етких продукционных сетей на~основе модели вывода Мамдани"--~Заде}
{Решение задачи классификации с использованием нейронных нечётких продукционных сетей \dots}

\ruauthor{О.~П.~Солдатова, И.~А.~Лёзин}
{О.~П.~С\,о\,л\,д\,а\,т\,о\,в\,а, И.~А.~Л\,ё\,з\,и\,н}

\ruaffil{\rusgau.}

\email{2}{\mem{op-soldatova@yandex.ru} (O.P.~Soldatova), \newline
\mem{ilyozin@yandex.ru} (I.~A.~Lyozin, \CAut)}

\ruauthordetails{2}{\fullauthor{\rupid{83143}{Ольга Петровна Солдатова}} \regalia{к.т.н., доц.}\position{доцент} \dept{каф.
информационных систем и технологий}
\fullauthor{\rupid{41321}{Илья Александрович Лёзин}} \regalia{к.т.н.}\position{доцент} \dept{каф.
информационных систем и технологий}
}

\rukeywords{задача классификации, нейронная нечёткая продукционная сеть,
многовыходовая сеть Ванга"--~Менделя, модель Мамдани"--~Заде}

\ruabstract{Рассматривается решение задачи распознавания объектов
пересекающихся классов с использованием систем нечеткого вывода и
нейронных сетей. Новая многовыходовая сеть Ванга"--~Менделя
сравнивается с новой архитектурой нейронной нечеткой продукционной
сети, основанной на модели Мамдани"--~Заде. Результаты
исследования данных моделей приведены при интерпретациях
логических операций, заданных соответственно алгебрами Гёделя,
Гогена и Лукашевича. Новая сеть Ванга"--~Менделя может
использовать минимум или основанную на сумме формулу как операции
$T$-нормы в соответствии с выбранной алгеброй вместо стандартной
операции произведения. Сеть Мамдани"--~Заде спроектирована в виде
каскада операций $T$-нормы, импликации и $S$-нормы, заданных
выбранной алгеброй. Кроме того, в сети Мамдани"--~Заде отсутствует
слой дефаззификации. Обе сети имеют несколько выходов в
соответствии с числом классов предметной области, что отличает их
от базовых реализаций. На выходах сетей формируются степени
принадлежности входного вектора заданным классам. Для сравнения
моделей использовались стандартные задачи классификации ирисов
Фишера и итальянских вин. В~данной статье приводятся результаты,
полученные при обучении сетей алгоритмом обратного распространения
ошибки. Анализ ошибок классификации показывает, что использование
данных алгебр в качестве интерпретации нечётких логических
операций, предложенное в статье, позволяет уменьшить погрешность
классификации как для многовыходовой сети Ванга"--~Менделя, так и
для новой сети Мамдани"--~Заде. Наилучшие результаты обучения
показывает алгебра Гёделя, но алгебра Лукашевича демонстрирует
лучшие обобщающие свойства при тестировании, что приводит к
наименьшему числу ошибок классификации.}

\entitle{Solving the Classification Problem by Using Neural Fuzzy Production Based Network Models of Mamdani--Zadeh}

\enauthor{O.~P.~Soldatova, I.~A.~Lyozin}{O.~P.~S\,o\,l\,d\,a\,t\,o\,v\,a, I.~A.~L\,y\,o\,z\,i\,n}

\enaffil{\ensgau.}

\enauthordetails{2}{\fullauthor{\enpid{83143}{Ol'ga P. Soldatova}} \regalia{Cand. Techn. Sci.}\position{Associate Professor} \dept{Dept. Information Systems and Technology}
\fullauthor{\enpid{41321}{Il'ya~A.~Lyozin}} \regalia{Cand. Techn. Sci.}\position{Associate Professor} \dept{Dept. Information Systems and Technology}
}



\enabstract{The article considers solving the problem of object recognition of intersected classes using fuzzy inference systems and neural networks. New multi-output network of Wang--Mendel is compared to a new architecture of neural fuzzy production network based on the model of Mamdani--Zadeh. Learning results of these models are given in the interpretation of logical operations provided by Godel, Goguen and Lukasiewicz algebras. New Wang--Mendel's network can use minimum or sum-based formula as $T$-norm operation in accordance with an appropriate algebra rather than the standard multiplication only. Mamdani--Zadeh's network is designed as a cascade of $T$-norm, implication and S-norm operations defined by selected algebra. Moreover defuzzification layer is not presented in Mamdani--Zadeh's network. Both networks have several outputs in accordance with the number of subject area classes what differs them from the basic realizations. Compliance degrees of an input vector to defined classes are formed at the network outputs. To compare the models the standard Fisher's irises and Italian wines classification problems were used. This article presents the results calculated by training the networks by backpropagation algorithm. Classification error analysis shows that the use of these algebras as interpreting fuzzy logic operations proposed in this paper can reduce the classification error for both multi-output network of Wang-Mendel and a new network of Mamdani--Zadeh. The best learning results are shown by Godel algebra, but Lukasiewicz algebra demonstrates better generalizing properties while testing, what leads to a less number of classification errors.}



 \enkeywords{classification problem, neural network
fuzzy production based network, multi-output network of
Wang-Mendel, model of Mamdani--Zadeh}

\date{03/X/2013}
\revisiondate{16/IV/2014}
\accepteddate{16/V/2014}

\makerutitle


При решении задачи распознавания объектов, принадлежащих пересекающимся
классам, успешно применяются нейронные нечёткие продукционные сети,
соединяющие возможности систем нечёткого вывода и нейронных сетей. 
Базовой моделью нечёткого вывода является модель Мамдани"--~Заде \cite{soldat:WM1, soldat:WM2}, однако широкое
распространение получила только одна нечёткая продукционная нейронная сеть "--- сеть Ванга"--~Менделя~[\citen{soldat:MZ1}--\citen{soldat:bib1}], которая была создана на её основе.

Применение сети Ванга"--~Менделя для решения задачи классификации вызывает
определённые трудности. Во"=первых, традиционная архитектура сети
предполагает один выходной нейрон, что приводит к увеличению погрешности
классификации при увеличении числа распознаваемых классов. Во"=вторых, в сети
Ванга"--~Менделя не предусмотрена модификация интерпретаций нечётких логических
операций, реализующих нечёткие продукционные правила вывода.

В~\cite{soldat:bib2} описана и исследована предложенная авторами многовыходовая модификация
сети Ванга"--~Менделя. В данной статье предлагается новая архитектура нейронной
нечёткой продукционной сети, основанной на модели Мамдани"--~Заде, и
исследуется решение задачи классификации для предложенной нечёткой сети и
сети Ванга"--~Менделя с несколькими выходами при интерпретациях нечётких
логических операций в соответствии в алгебрами Гёделя, Гогена и Лукашевича.

Типовая структура системы нечёткого вывода включает в себя блок
фаззификации, базу нечётких правил вывода, модуль вывода решения и блок
дефаззификации.

Фаззификатор преобразует чёткое множество входных данных \linebreak $x=(x_1,x_2
,\ldots,x_N)^\top\in X$ в нечёткое множество $A\subseteq X$, определённое с помощью конкретной функции
принадлежности $\mu _A ( x )$; здесь $$A=A_1 \times A_2 \times \ldots \times A_N.$$ 
Дефаззификатор преобразует несколько нечётких множеств $B^k$, $k=1, 2, \ldots,$ $M$, или единственное нечёткое
множество $B$ в конкретное значение выходной переменной $y \in Y$ на основании нечётких выводов, вырабатываемых модулем вывода решений, и в соответствии с выбранным методом дефаззификации (символами $X$ и $Y$ обозначаются пространства входных и выходных переменных).

Подобная система также может быть представлена в несколько ином виде, в частности на рис.~\ref{lez:pict1}~\cite{soldat:bib1}.

\begin{figure}[b!]
\centering \footnotesize
\includegraphics[scale=0.86]{soldat1}
\caption{Структура системы нечёткого вывода при $M$ правилах вывода \label{lez:pict1}}

\smallskip

[Figure \ref{lez:pict1}. Architecture of fuzzy output by using $M$ rules]

\end{figure}

Выходной сигнал модуля вывода решений может иметь вид $M$ нечётких
множеств, определяющих диапазон изменения выходной переменной. Агрегатор
преобразует $M$ нечётких множеств в одно нечёткое множество, а
дефаззификатор преобразует это множество в одно конкретное значение,
принимаемое в качестве выходного сигнала всей системы.

В модели вывода Мамдани"--~Заде присутствуют следующие операции~\cite{soldat:bib1}:
\begin{itemize}
\item операция логического или арифметического произведения для определения значения функции принадлежности условий правил, в которой учитываются все компоненты вектора условия;
\item операция логического или арифметического произведения для определения значения функции принадлежности для всей импликации ${A^k {\to} B^k}$;
\item операция логической суммы для агрегации равнозначных результатов импликации многих правил;
\item оператор дефаззификации, трансформирующий нечёткий результат \linebreak $\mu _{B'} (y)$ в чёткое значение $y$.
\end{itemize}

Считается, что все $M$ правил связаны между собой логической операцией
<<ИЛИ>>, а выходы правил $y^1, y^2,\ldots,y^M$ взаимно независимы.
Следовательно, можно использовать правила вида:
\begin{equation}
\label{soldat:eq1}
\!\!
\begin{array}{l}
R^1 : \text{ IF }  x_1^1  \text{ is } A_1^1   \text{ AND } x_2^1 \text{ is }  A_2^1 \text{ AND } \dots  \text{ AND }
 x_N^1  \text{ is }  A_N , \text{ THEN } y^1  \text{ is } B^1;\\
R^k :  \text{ IF }   x_1^k   \text{ is } A_1^k   \text{ AND }  x_2^k   \text{ is } A_2^k
  \text{ AND } \dots   \text{ AND }  x_N^k   \text{ is } A_N , \text{ THEN } y^k \text{ is } B^k;\\
R^M :  \text{ IF }  x_1^M  \text{ is } A_1^M  \text{ AND }  x_2^M  \text{ is } A_2^M
 \text{ AND } \dots  \text{ AND }  x_N^M   \text{ is } A_N , \\
\hspace{9.5cm} 
  \text{ THEN } y^M \text{ is } B^M.
\end{array}\!\!
\end{equation}

Переменные $(x_1^k ,x_2^k ,\ldots,x_N^k )^\top$ образуют $N$-мерный входной вектор $x^k$, составляющий аргумент условия, в котором $A_1^k$, $A_2^k,\ldots,A_N^k$ и $B^k$ обозначают нечёткие множества, а $X_1$, $X_2,\ldots,X_N,Y$ "--- чёткие
множества, причём $A_1^k \subseteq X_1$, $A_2^k \subseteq X_2$, $\ldots$, $A_N^k \subseteq X_N $, а $B^k\subseteq Y$. 
Если обозначить декартово произведение
множеств $A_1^k$, $A_2^k ,\ldots,A_N^k$ как $A^k=A_1^k \times A_2^k
\times \ldots \times A_N^k$, то соответствующие значения коэффициентов
принадлежности условия и заключения правила вывода "--- $\mu_{A^k} (x^k)$ и $\mu_{B^k} (y^k)$ соответственно. 
Тогда правило типа \eqref{soldat:eq1} можно представить в виде нечёткой импликации:
\begin{equation}
\label{soldat:eq2}
R^k: \quad A^k\to B^k,\quad k=1, 2, \ldots,M.
\end{equation}

В то же время правило типа \eqref{soldat:eq1} может быть интерпретировано как нечёткое
отношение, определённое на множестве $X\times Y$, где $X$ "--- декартово
произведение чётких множеств и $X=X_1 \times X_2 \times \ldots \times X_N$,
причём $\left(x_1^k,x_2^k ,\ldots,x_N^k \right)\in X$, $x_1^k \in X_1$, $x_2^k \in X_2
$, $\ldots$, $x_N^k \in X_N $, а $y^k\in Y$, то есть $R^k\subseteq X\times Y$ "--- 
нечёткое множество с функцией принадлежности~\cite{soldat:bib3}:
\begin{equation}
\label{soldat:eq3}
\mu_{R^k}(x^k,y^k)=\mu_{A^k\to B^k} (x^k,y^k).
\end{equation}

Если на выходе модуля вывода решения получается $M$ нечётких множеств $B^k\subseteq Y$, то нечёткое множество $B^k$ с функцией принадлежности
\begin{equation}
\label{soldat:eq4}
\mu_{B^k} (y^k)=\sup \limits_{x\in X} 
\left[\mu_{A^k} (x^k) \genfrac{}{}{0pt}{1}{T}{\ast}\mu _{A^k\to B^k} (x^k,y^k)\right]
\end{equation}
определяется как комбинация нечёткого множества $A^k$ и отношения $R^k$:
\begin{equation}
\label{soldat:eq5}
B^k=A^k\circ R^k,\quad k=1, 2, \dots, M.
\end{equation}
Здесь $\genfrac{}{}{0pt}{1}{T}{\ast}$ "--- обобщённое понятие оператора $T$-нормы \cite{soldat:bib3}.
Таким образом, конкретная форма $\mu _{B^k} (y^k)$ зависит от
применяемой $T$-нормы, определения нечёткой импликации $R^k$ и от способа
определения декартова произведения нечётких множеств~\cite{soldat:bib3}.

Если на выходе модуля вывода решения получается одно нечёткое множество
$B\subseteq Y$ (в системе есть агрегатор), то 
$$B=A\circ
\bigcup \limits_{k=1}^M R^k =A\circ R,$$ где 
$$A=A_1 \times A_2 \times \ldots \times A_N , \quad R=\bigcup\limits_{k=1}^M {R^k} .$$ 
В этом случае
\begin{equation}
\label{soldat:eq6}
\mu _B ( y )=\sup \limits_{x\in X} \Bigl[ \mu _A
(x) \genfrac{}{}{0pt}{1}{T}{\ast}
\Bigl(
\mathop{\genfrac{}{}{0pt}{1}{S}{\ast}}
\limits_{1\le k\le M}
\bigl(\mu _{R^k}
(x^k,y)\bigr) \Bigr)\Bigr].
\end{equation}
Здесь $\mathop{\genfrac{}{}{0pt}{1}{S}{\ast}}
\limits_{1\le k\le M} ({}\cdot{})$ "--- обобщённое понятие оператора $S$-нормы над вектором операндов в круглых скобках~\cite{soldat:bib3}.
Таким образом, конкретная форма $\mu _B ( y )$ зависит ещё и от применяемой в агрегаторе операции $S$-нормы.

Так как доказано, что любая непрерывная $T$-норма изоморфна либо произведению,
либо конъюнкции Лукашевича, либо эквивалентна минимуму, или в смысле
порядковой суммы является их <<смесью>>, то основными $T$-нормами являются:
минимум, произведение и конъюнкция Лукашевича. Данный результат упрощает
выбор интерпретации нечётких логических операций. Двойственными для данных
$T$-норм соответственно являются следующие $S$-нормы: максимум, вероятностная
сумма и дизъюнкция Лукашевича. При этом операциями деления, соответствующим
базовым $T$-нормам, являются импликация Гёделя, импликация Гогена и импликация
Лукашевича~\cite{soldat:bib4}.

Таким образом, в модели нечёткого вывода Мамдани"--~Заде имеет смысл
рассматривать интерпретации логических операций, заданныx соответственно
алгебрами Гёделя, Гогена и Лукашевича. Если обозначить $T$-норму как <<$\otimes
$>>, $S$-норму как <<$\oplus $>>, а импликацию как <<$\to $>>, то нижеследующие формулы~\eqref{soldat:eq7}
определяют операции в соответствии с алгеброй Гёделя (Godel's algebra), \eqref{soldat:eq8} "--- в соответствии
с алгеброй Гогена (Goguen's algebra), а формулы \eqref{soldat:eq9} "--- в~соответствии с алгеброй Лукашевича (Lukasiewicz's algebra):
\begin{gather}
\label{soldat:eq7}
\begin{cases}
\mu _A ( x )\otimes \mu _B (y) & = \min \left\{ \mu _A (x),\mu _B (y) \right\}; \\
\mu _{A\to B} ( x,y) &= \begin{cases}
~~~~1, & \text{если} \quad \mu _A ( x )\le \mu _B (y),\\
\mu _B (y), & \text{если} \quad \mu _A ( x )>\mu _B (y); \\
\end{cases} \\
\mu _A ( x )\oplus \mu _B (y) & = \max \left\{ \mu _A ( x
),\mu _B (y)\right\};
\end{cases}\\
\label{soldat:eq8}
\begin{cases}
\mu _A ( x )\otimes \mu _B (y) &= \mu _A ( x )  \mu _B (y); \\
\mu _{A\to B} ( x,y)& =\begin{cases}
~~~~~~1,\quad &\text{если} \quad \mu _A ( x )\le \mu_B (y), \\
{\mu _B (y)}/{\mu _A ( x )},\quad &\text{если} \quad \mu _A ( x )>\mu _B (y);
\end{cases} \\
\mu _A ( x )\oplus \mu _B (y) & = \mu _A ( x )+\mu _B
(y)-\mu _A ( x )  \mu _B (y);
\end{cases}\\
\label{soldat:eq9}
\begin{cases}
\mu _A ( x )\otimes \mu _B (y) & = \max \{ 0, \mu _A ( x )+\mu _B (y)-1 \}; \\[2mm]
\mu _{A\to B} (x,y) &=\min \{ 1, 1-\mu _A ( x )+\mu _B (y) \}; \\[2mm]
\mu _A ( x )\oplus \mu _B (y) &=\min \{ 1, \mu _A ( x )+\mu _B (y) \}.
\end{cases}
\end{gather}

Системы нечёткого вывода можно представить в виде многослойных нейронных
сетей с прямым распространением сигнала, которые получили название нейронных
нечётких продукционных сетей~\cite{soldat:bib5}. Для обучения подобных сетей, как правило,
используются градиентные алгоритмы и метод обратного распространения ошибки.
Условие дифференцируемости функций активации нейронов обусловило применение
в подобных моделях следующих интерпретаций нечётких логических операций:

\begin{itemize}
\item в качестве операции $T$-нормы и импликации используется алгебраическое произведение;
\item в качестве декартова произведения используется алгебраическое произведение или минимум;
\item аккумулирование заключений правил вывода не проводится или в качестве операции $S$-нормы используется алгебраическая или взвешенная сумма;
\item в качестве функций фаззификации используется функция Гаусса или сигмоидальные функции;
\item дефаззификация выполняется по методу центра тяжести или по методу среднего центра.
\end{itemize}

В частности, в сети Ванга"--~Менделя "--- самой известной реализации модели
нечёткого вывода Мамдани"--~Заде используются следующие ограничения~\cite{soldat:bib5}:
\begin{itemize}
\item входные переменные являются чёткими;
\item функции фаззификации "--- функции Гаусса;
\item декартово произведение "--- в форме минимума;
\item нечёткая импликация "--- произведение;
\item $T$-норма "--- произведение;
\item аккумулирование заключений правил не проводится;
\item метод дефаззификации "--- метод среднего центра.
\end{itemize}

В нейроимитаторе <<Нейрокомбайн>>, описанном в~\cite{soldat:bib2}, предложена и реализована
другая архитектура сети Ванга"--~Менделя, в которой предусмотрено несколько
выходных нейронов. Структура сети с двумя входами, тремя правилами вывода и
двумя выходами приведена на рис.~\ref{lez:pict2}. В отличие от стандартной
четырёхслойной конфигурации сети Ванга"--~Менделя, данная модель состоит из
пяти слоёв нейронов.

\begin{figure}[b!]
\centering \footnotesize
\includegraphics[scale=0.9]{soldat2}
\caption{Структура нейронной нечёткой сети Ванга"--~Менделя с двумя выходами \label{lez:pict2}}

\smallskip

[Figure \ref{lez:pict2}. Architecture of Wang--Mendel's neural fuzzy network with two outputs]

\vspace{-3mm}
\end{figure}


Первый и второй слои предложенной модели не отличаются от стандартной сети
Ванга"--~Менделя и реализуют фаззификацию по функции Гаусса $N$-мерного входного
вектора $x$, а также агрегацию условий правил вывода в соответствии с
операцией $T$-нормы в форме произведения. Величины $\mu_A^k (x_j)$ задают функции принадлежности входных переменных $x_j$ к нечёткому множеству $A$ в соответствии с правилом $k$. Величины $w_k$ обозначают
степень принадлежности условия $k$-того правила, полученную в результате
агрегации.

Третий слой реализует операцию импликации в форме произведения (это
параметрический слой). В процессе обучения подбираются параметры $v_{ks}$,
соответствующие степени принадлежности $\mu _{B^k} ( y_s^k )$,
где $k$ "--- номер правила вывода, а $s$ "--- индекс выходного нейрона,
соответствующий номеру класса, к которому принадлежит входной вектор.
Величины $v_{ks}$ соответствуют весам связей нейронов третьего и четвёртого
слоёв. Величины $z_{ks} $ означают результаты операции нечёткой импликации,
которая интерпретируется в данной модели в форме произведения.

Четвёртый слой осуществляет агрегирование $M$ правил вывода (первый и второй
нейрон) и генерацию нормализующего сигнала (третий нейрон). Результаты
агрегации обозначены для первого и второго нейрона как $f_{11} $ и $f_{12}
$ соответственно. Нормализующие сигналы обозначены как $f_{21} $ и $f_{22}
$ соответственно. Число агрегирующих нейронов в данном слое равно числу
выходов сети. В отличие от стандартной модели "--- непараметрический слой.

Пятый слой состоит из двух выходных нейронов и выполняет нормализацию,
формируя выходной сигнал $y_s $ для каждого класса. Это непараметрический
слой.

На рис.~\ref{lez:pict2} знаки $\times$, $+$, $/$ соответствуют алгебраическим операциям умножения, сложения и деления.


Предложенная структура нейронной сети легко модифицируется на случай с~числом выходов, большим, чем два. Таким образом, нейронная сеть реализует функции, которые можно записать в виде
\begin{multline}
\label{soldat:eq10}
y_s (x)=\frac{1}{\sum_{k=1}^M {\left[ \prod_{j=1}^N \mu _A^k
(x_j ) \right]} }\cdot \biggl( \sum\limits_{k=1}^M v_{ks} \cdot \biggl[
\prod\limits_{j=1}^N \mu _A^k (x_j ) \biggr]
\biggr)=
\\
=
\frac{1}{\sum_{k=1}^M w_k  }\cdot \biggl(
\sum\limits_{k=1}^M z_{ks}  \biggr).
\end{multline}

Слой дефаззификации в приведённой выше модели нейронной сети позволяет
сформировать на выходе сети чёткое значение, что необходимо при решении
задачи прогнозирования или задачи аппроксимации функций. Для решения задач
классификации, идентификации или распознавания чёткое значение на выходе
сети не является обязательным, а при пересекающихся классах объектов вообще
не имеет смысла, так как на выходе сети требуется получить степень
принадлежности предъявленного входного вектора к конкретному классу. В этой
связи авторами предлагается новая модификация модели нейронной нечёткой
продукционной сети Мамдани"--~Заде, исследованная на примере решения задачи
классификации.

Структура сети с двумя входами, четырьмя правилами вывода и двумя выходами
представлена на рис.~\ref{lez:pict3}.

\begin{figure}[h!]
\centering \footnotesize
\includegraphics{soldat3}
\caption{Структура нейронной нечёткой продукционной сети Мамдани"--~Заде двумя выходами \label{lez:pict3}
[Figure \ref{lez:pict3}. Architecture of Mamdani--Zadeh's neural fuzzy network with two outputs]}

\vspace{-5mm}

\end{figure}


Для данной сети используются следующие ограничения:
\begin{itemize}
\item входные переменные являются чёткими;
\item функции фаззификации "--- функции Гаусса;
\item декартово произведение "--- в форме минимума;
\item нечёткая импликация "--- произведение, минимум или импликация Лукашевича;
\item $T$-норма "--- произведение, минимум или конъюнкция Лукашевича;
\item аккумулирование заключений правил проводится в соответствии с операцией $S$-нормы "--- вероятностной суммы, максимума или дизъюнкции Лукашевича;
\item дефаззификация не проводится.
\end{itemize}

Первый и второй слои предложенной модели не отличаются по функциональности
от модели, представленной на рис.~\ref{lez:pict2}, за исключением знака крупного
$\times$ для обозначения операции $T$-нормы, которая имеет приведённые выше
интерпретации.

Третий слой реализует операцию нечёткой импликации (знак $\to$) в
соответствии с приведёнными выше интерпретациями (это параметрический слой).
В процессе обучения подбираются параметры $v_k$, соответствующие степени
принадлежности $\mu_{B^k} ( y^k)$, где $k$ "---  номер
правила вывода. Величины $z_k$ означают результаты операции нечёткой
импликации.

Четвёртый слой осуществляет агрегирование $M$ правил вывода (первый и второй
нейрон) для каждого класса входных векторов в соответствии с интерпретациями
нечёткой дизъюнкции (знак крупный $+$) и формирует выходной сигнал $y_s$ для
каждого класса. Это непараметрический слой.

Исследования проводились на данных репозитория UCI (Machine Learning
Repository; {\tt \url{http://archive.ics.uci.edu/ml}}), который представляет собой набор реальных и модельных задач
машинного обучения, используемых для эмпирического анализа алгоритмов
машинного обучения. 
Репозиторий содержит реальные данные по прикладным
задачам, которые применяются для оценки точности классификаторов, в том
числе и на основе нечётких продукционных баз знаний и нейронных сетей~\cite{soldat:bib7,soldat:bib8}.

В настоящей работе использовались данные для решения задач классификации ирисов
Фишера и итальянских вин. В исследовании сравниваются нейронные сети,
структуры которых представлены на рис.~\ref{lez:pict2} и~\ref{lez:pict3}.

Ирисы Фишера "--- набор данных, на примере которого Рональд Фишер в~1936 году продемонстрировал разработанный им метод дискриминантного анализа. 
Ирисы Фишера состоят из данных о~150 экземплярах ирисов, по 50~экземпляров трёх видов "--- ирис щетинистый (англ. \textit{Iris setosa}), ирис виргинский (англ. \textit{Iris virginica}) и ирис разноцветный (англ. \textit{Iris versicolor}). 
Для каждого экземпляра измерялись четыре характеристики. Один из классов (\textit{Iris setosa}) линейно отделим от двух остальных. 
Для обучения использовались 90 образцов ($\approx 60$\,{\%}), а оставшиеся 60~образцов использовались для тестирования качества решения задачи.

Набор данных задачи классификации итальянских вин представляет собой результаты химического анализа вин, принадлежащих трём различным сортам. 
В~ходе анализа было выделено процентное содержание 13 признаков, присутствующих в каждом из трёх сортов вин. Общий объём данных "--- 178 образцов. 
Из них на 142 ($\approx 80$\,{\%}) образцах проводилось обучение, а на 36 "--- тестирование качества решения задачи. Все классы пересекаются.

В ходе экспериментальных исследований при решении задачи классификации
ирисов была определена зависимость относительной погрешности обучения от
количества правил и используемой алгебры, а также относительная погрешность
классификации при следующих параметрах сети: число нейронов во входном слое
"--- 4, число нейронов в выходном слое "--- 3, число итераций обучения "--- 2500.
Результаты исследования представлены в табл. \ref{lez:tabl1} и \ref{lez:tabl2}.

В качестве алгоритма обучения использовался алгоритм обратного
распространения ошибки. Ввиду недифференцируемости операций минимума и
максимума в интерпретациях нечётких логических операций конъюнкции и
дизъюнкции в алгебре Гёделя и нечётких логических операций конъюнкции,
дизъюнкции и импликации в алгебре Лукашевича был применён метод
корректировки весов сети только для синаптических связей, соответствующих
минимуму или максимуму значений, описанный в~\cite{soldat:bib3}.

Однако указанный подход не может быть применён для операций нечёткой
импликации в алгебрах Гёделя и Гогена, вследствие чего обучение сети
Мамдани"--~Заде с алгебрами Гёделя и Гогена реализовать невозможно.

\begin{sidewaystable}
\centering
\begin{minipage}{.94\textwidth}
\centering\small
\caption{\bf Зависимость относительной погрешности обучения от количества правил
и~используемой алгебры для задачи классификации ирисов [Dependence of the learning relative error on the number of rules and algebra used for classification of the irises] \label{lez:tabl1} }\vspace{-5mm}
\begin{tabular}{c|c|c|c|c} \hline
\raisebox{-4ex}[0cm][0cm]{\backslashbox{\footnotesize Правило [Rule]}{\footnotesize Алгебра [Algebra]}}& \multicolumn{3}{c|}{\footnotesize Многовыходовая сеть Ванга"--~Менделя} &\rule{0mm}{12pt}\footnotesize Сеть Мамдани"--~Заде \\ 
& \multicolumn{3}{c|}{\footnotesize [Multi-output Wang--Mendel's networks]} & \footnotesize [Mamdani--Zadeh's network] \\
\cline{2-5}
 &\rule{0mm}{12pt}\footnotesize Алгебра Гёделя  &\footnotesize Алгебра Гогена  &\footnotesize Алгебра  Лукашевича    &\footnotesize Алгебра  Лукашевича  \\
 &\footnotesize [Godel's algebra]  &\footnotesize [Goguen's algebra]    &\footnotesize [Lukasiewicz's algebra] &\footnotesize [Lukasiewicz's algebra] \\ \hline
\rule{0mm}{12pt}3& 0.056& 0.044 & \textbf{0.011} & \textbf{0.044} \\ 
4& 0.056& 0.033 & 0.044          & 0.056 \\ 
5& 0.056& 0.022 & 0.044          & 0.056 \\ 
6& \textbf{0.044} & 0.044 & 0.022 & \textbf{0.044} \\ 
7& 0.056& 0.022 & 0.033          & \textbf{0.044} \\ 
8& 0.056& \textbf{0.011}& 0.044  & \textbf{0.044} \\ 
9& \textbf{0.044} & 0.044 & 0.044& \textbf{0.044} \\ \hline
\end{tabular}


\smallskip

\bigskip

\caption{\bf Зависимость относительной погрешности классификации от~количества
правил и~используемой алгебры для~задачи классификации ирисов при~проверке
на тестовой выборке [Dependence of the  classification relative error on the number of rules and algebra used for classification of the irises on the testing set]\label{lez:tabl2}}\vspace{-3mm}
\begin{tabular}{c|c|c|c|c} \hline
\raisebox{-4ex}[0cm][0cm]{\backslashbox{\footnotesize Правило [Rule]}{\footnotesize Алгебра [Algebra]}}& \multicolumn{3}{c|}{\footnotesize Многовыходовая сеть Ванга"--~Менделя} &\rule{0mm}{12pt}\footnotesize Сеть Мамдани"--~Заде \\ 
& \multicolumn{3}{c|}{\footnotesize [Multi-output Wang--Mendel's networks]} & \footnotesize [Mamdani--Zadeh's network] \\
\cline{2-5}
 &\rule{0mm}{12pt}\footnotesize Алгебра Гёделя  &\footnotesize Алгебра Гогена  &\footnotesize Алгебра  Лукашевича    &\footnotesize Алгебра  Лукашевича  \\
 &\footnotesize [Godel's algebra]  &\footnotesize [Goguen's algebra]    &\footnotesize [Lukasiewicz's algebra] &\footnotesize [Lukasiewicz's algebra] \\ \hline
\rule{0mm}{12pt}3& \textbf{0.082}& 0.061& 0.082& \textbf{0.041} \\ 
4& 0,102 & 0.061 & 0.061& 0.061 \\ 
5& \textbf{0.082}& \textbf{0.041} & 0,102 & \textbf{0.041} \\ 
6& \textbf{0.082}& \textbf{0.041} & 0.041 & \textbf{0.041} \\ 
7& \textbf{0.082}& 0.082 & \textbf{0.020}  & \textbf{0.041} \\ 
8& \textbf{0.082}& \textbf{0.041} & 0.041 & 0.061 \\ 
9& \textbf{0.082}& \textbf{0.041} & 0.061 & \textbf{0.041} \\ \hline
\end{tabular}

\end{minipage}

\end{sidewaystable}



\begin{sidewaystable}
\centering
\begin{minipage}{.94\textwidth}
\centering\small
\caption{\bf Зависимость относительной погрешности обучения от~количества правил
и~используемой алгебры для~задачи классификации вин [Dependence of the learning relative error on the number of rules and algebra used for classification of the wines] \label{lez:tabl3} }\vspace{-5mm}
\begin{tabular}{c|c|c|c|c} \hline
\raisebox{-4ex}[0cm][0cm]{\backslashbox{\footnotesize Правило [Rule]}{\footnotesize Алгебра [Algebra]}}& \multicolumn{3}{c|}{\footnotesize Многовыходовая сеть Ванга"--~Менделя} &\rule{0mm}{12pt}\footnotesize Сеть Мамдани"--~Заде \\ 
& \multicolumn{3}{c|}{\footnotesize [Multi-output Wang--Mendel's networks]} & \footnotesize [Mamdani--Zadeh's network] \\
\cline{2-5}
 &\rule{0mm}{12pt}\footnotesize Алгебра Гёделя  &\footnotesize Алгебра Гогена  &\footnotesize Алгебра  Лукашевича    &\footnotesize Алгебра  Лукашевича  \\
 &\footnotesize [Godel's algebra]  &\footnotesize [Goguen's algebra]    &\footnotesize [Lukasiewicz's algebra] &\footnotesize [Lukasiewicz's algebra] \\ \hline
\rule{0mm}{12pt}3& 0.056 & \textbf{0} & 0.021 & 0.028 \\ 
4& 0.035 & \textbf{0} & 0.028 & 0.014 \\  
5& 0.049 & \textbf{0} & \textbf{0.007} & 0.021 \\ 
6& 0.063 & \textbf{0} & 0.028 & 0.028 \\ 
7& \textbf{0.028} & \textbf{0}& 0.021 & \textbf{0.007} \\ 
8& 0.049 & \textbf{0} & 0.028 & 0.021 \\ 
9& \textbf{0.028}& \textbf{0} & 0.021 & 0.021 \\ \hline
\end{tabular}

\smallskip
\bigskip

\caption{\bf Зависимость относительной погрешности классификации от~количества
правил и~используемой алгебры для~задачи классификации вин при~проверке на~тестовой выборке [Dependence of the  classification relative error on the number of rules and algebra used for classification of the wines on the testing set]\label{lez:tabl4} }\vspace{-3mm}
\begin{tabular}{c|c|c|c|c} \hline
\raisebox{-4ex}[0cm][0cm]{\backslashbox{\footnotesize Правило [Rule]}{\footnotesize Алгебра [Algebra]}}& \multicolumn{3}{c|}{\footnotesize Многовыходовая сеть Ванга"--~Менделя} &\rule{0mm}{12pt}\footnotesize Сеть Мамдани"--~Заде \\ 
& \multicolumn{3}{c|}{\footnotesize [Multi-output Wang--Mendel's networks]} & \footnotesize [Mamdani--Zadeh's network] \\
\cline{2-5}
 &\rule{0mm}{12pt}\footnotesize Алгебра Гёделя  &\footnotesize Алгебра Гогена  &\footnotesize Алгебра  Лукашевича    &\footnotesize Алгебра  Лукашевича  \\
 &\footnotesize [Godel's algebra]  &\footnotesize [Goguen's algebra]    &\footnotesize [Lukasiewicz's algebra] &\footnotesize [Lukasiewicz's algebra] \\ \hline
\rule{0mm}{12pt}3& 0,111 & \textbf{0.028} & 0.056 & 0.083 \\ 
4& \textbf{0.028} & \textbf{0.028}& 0.056 & 0.056 \\ 
5& 0.056 & 0.056 & 0.083 & 0.083 \\ 
6& 0.056 & 0.056 & 0.111 & 0.056 \\ 
7& 0.056 & 0.056 & 0.083 & 0.083 \\ 
8& 0.056 & 0.056 & \textbf{0.028} & 0.056 \\ 
9& 0.056 & 0.083 & 0.111 & \textbf{0.028} \\ \hline
\end{tabular}
\end{minipage}

\end{sidewaystable}


Как видно из таблиц, в среднем лучшие результаты обучения показывает сеть Ванга"--~Менделя с использованием алгебры Гогена или Лукашевича, но при этом более ровную погрешность классификации на тестовой выборке демонстрирует сеть Мамдани"--~Заде, что характеризует лучшую сходимость обучения.

В ходе экспериментальных исследований при решении задачи классификации вин была определена зависимость относительной погрешности обучения от количества правил и используемой алгебры, а также относительная погрешность
классификации при следующих параметрах сети: число нейронов во входном слое
"--- 13, число нейронов в выходном слое "--- 3, число итераций обучения "--- 2500.
Результаты исследования представлены в табл.~\ref{lez:tabl3} и~\ref{lez:tabl4}.

Как видно из таблиц, лучшие результаты обучения демонстрирует сеть
Ванга"--~Менделя, использующая алгебру Гогена, но при этом обе сети на тестовых
выборках показывают в среднем аналогичные результаты с использованием
алгебры Лукашевича, что согласуется с теоретическими исследованиями,
представленными в~\cite{soldat:bib4}. Алгебра Лукашевича в модифицированной сети
Мамдани"--~Заде, в частности, проявляет большую способность к выделению общих
признаков, а не адаптацию к обучающей выборке, что характеризуется более
точной работой на тестовых выборках. Более высокая погрешность сети
Ванга"--~Менделя объясняется тем, что в данной сети можно модифицировать только
операцию нечёткой конъюнкции, в то время как операции нечёткой импликации и
нечёткой дизъюнкции не используются вообще. В~свою очередь, модифицированная
сеть применяет интерпретации всех трёх операций.


\begin{thebibliography}{99}

\Bibitem{soldat:WM1}
\by L.~X.~Wang, J.~M.~Mendel
\crossref{10.1109/isic.1991.187368}
\paper Generating fuzzy rules by learning from examples
\jour IEEE Trans. Syst., Man, Cybern.
\vol 22
\issue 6 
\yr 1992
\pages 1414--1427


\Bibitem{soldat:WM2}
\paper The WM method completed: a flexible fuzzy system approach to data mining
\jour IEEE Trans. Fuzzy Systems 
\vol 11 
\issue 6 
\yr 2003
\pages 768--782
\by Li-Xin~Wang
\crossref{10.1109/TFUZZ.2003.819839}

\Bibitem{soldat:MZ1}
\by L.~A.~Zadeh
\paper Fuzzy logic, neural networks, and soft computing
\jour Communications of the ACM
\yr 1994
\vol 37
\issue 3
\pages 77--84

\Bibitem{soldat:MZ2}
\paper  Application of Fuzzy Logic to Approximate Reasoning Using Linguistic Synthesis
\jour IEEE Trans. Computers 
\vol C-26 
\issue 12 
\pages 1182--1191
\by E.~H.~Mamdani
\crossref{10.1109/tc.1977.1674779}


\RBibitem{soldat:bib1}
\by С.~Осовский
\publaddr М.
\book Нейронные сети для обработки информации
\publ Финансы и статистика
\yr 2002
\totalpages 344
\misctransl
\lang In Russian
\by S.~Osovskiy
\book Neyronnyye seti dlya obrabotki informatsii \rm [Neural networks for information processing]
\publ Finansy i statistika
\yr 2002
\totalpages 344
\publaddr Moscow


\RBibitem{soldat:bib2}
\by О.~П.~Солдатова
\paper Многофункциональный имитатор нейронных сетей
\jour Программные продукты и системы
\yr 2012
\issue 3
\pages 27--31
\elib{http://elibrary.ru/item.asp?id=19685733}
\misctransl
\lang In Russian
\paper Multifunctional simulator of neural networks 
\by O.~P.~Soldatova
\jour Programmnyye produkty i sistemy
\yr 2012
\issue 3
\pages 27--31

\RBibitem{soldat:bib3}
\by Д.~Рутковская, М.~Пилиньский, Л.~Рутковский
\book Нейронные сети, генетические алгоритмы и нечёткие системы
\publaddr М.
\publ Горячая линия--Телеком
\yr 2007
\totalpages 452
\misctransl
\lang In Russian
\by D.~Rutkovskaya, M.~Pilin'skiy, L.~Rutkovskiy
\book Neyronnyye seti, geneticheskiye algoritmy i nechotkiye sistemy \rm [Neural networks, genetic algorithms and fuzzy systems]
\publaddr M.
\publ Goryachaya liniya--Telekom
\yr 2007
\totalpages 452

\Bibitem{soldat:bib4}
\book Mathematical Principles of Fuzzy Logic
\serial The Springer International Series in Engineering and Computer Science
\vol 517
\publ Springer
\by V.~Nov\'ak, I.~Perfilieva, J.~Mo\v{c}ko\v{r}
\yr 1999
\totalpages xiii+320
\crossref{10.1007/978-1-4615-5217-8}
\rtransl 
\by В.~Новак, И.~Перфильева, И.~Мочкорж
\book Математические принципы нечёткой логики
\publ М.
\publaddr Физматлит
\yr  2006
\totalpages 352

\RBibitem{soldat:bib5}
\by В.~В.~Борисов, В.~В.~Круглов, А.~С.~Федулов
\book Нечеткие модели и сети
\publaddr М.
\publ Горячая линия--Телеком
\yr 2007
\totalpages 284
\misctransl
\by V.~V.~Borisov, V.~V.~Kruglov, A.~S.~Fedulov
\book Nechetkiye modeli i seti \rm [Fuzzy models and networks]
\publaddr Moscow
\publ Goryachaya liniya--Telekom
\yr 2007
\totalpages 284


\RBibitem{soldat:bib7}
\by А.~С.~Катасёв
\paper Математическое обеспечение и программный комплекс формирования нечётко-продукционных баз знаний для экспертных диагностических систем
\jour Фундаментальные исследования
\yr 2013
\issue 10-9
\pages 1922--1927
\elib{http://elibrary.ru/item.asp?id=20683012}
\misctransl
\lang In Russian
\paper Mathematical and software for fuzzy-productions knowledge bases generation of the expert diagnostic systems
\by A.~S.~Katasev
\jour Fundamental'nyye issledovaniya
\yr 2013
\issue 10-9
\pages 1922--1927 


\RBibitem{soldat:bib8}
\by В.~В.~Бухтояров
\paper Трехступенчатый эволюционный метод формирования коллективов нейронных сетей для решения задач классификации
\jour Программные продукты и системы
\yr 2012
\issue 4
\pages 101--106
\elib{http://elibrary.ru/item.asp?id=19685830}
\misctransl
\lang In Russian
\by V.~V.~Bukhtoyarov
\paper Evolutionary three-stage approach for designing of neural networks ensembles for classification problems
\jour Programmnyye produkty i sistemy
\yr 2012
\issue 4
\pages 101--106


\end{thebibliography}


\renewcommand{\correctvfill}{\vspace{-13mm}}


%
%\chead[\fancyplain{}{\scriptsize \sl
%O.~P.~S\,o\,l\,d\,a\,t\,o\,v\,a, I.~A.~L\,y\,o\,z\,i\,n}]{\fancyplain{}{\scriptsize
%\sl  Solving the Classification Problem by Using Neural Fuzzy Production \dots}}

\makeentitle

\renewcommand{\baselinestretch}{0.9}

 \newpage
%
% \tableofengcontents
%
%\end{document}
